% Options for packages loaded elsewhere
\PassOptionsToPackage{unicode}{hyperref}
\PassOptionsToPackage{hyphens}{url}
%
\documentclass[
]{article}
\usepackage{amsmath,amssymb}
\usepackage{iftex}
\ifPDFTeX
  \usepackage[T1]{fontenc}
  \usepackage[utf8]{inputenc}
  \usepackage{textcomp} % provide euro and other symbols
\else % if luatex or xetex
  \usepackage{unicode-math} % this also loads fontspec
  \defaultfontfeatures{Scale=MatchLowercase}
  \defaultfontfeatures[\rmfamily]{Ligatures=TeX,Scale=1}
\fi
\usepackage{lmodern}
\ifPDFTeX\else
  % xetex/luatex font selection
\fi
% Use upquote if available, for straight quotes in verbatim environments
\IfFileExists{upquote.sty}{\usepackage{upquote}}{}
\IfFileExists{microtype.sty}{% use microtype if available
  \usepackage[]{microtype}
  \UseMicrotypeSet[protrusion]{basicmath} % disable protrusion for tt fonts
}{}
\makeatletter
\@ifundefined{KOMAClassName}{% if non-KOMA class
  \IfFileExists{parskip.sty}{%
    \usepackage{parskip}
  }{% else
    \setlength{\parindent}{0pt}
    \setlength{\parskip}{6pt plus 2pt minus 1pt}}
}{% if KOMA class
  \KOMAoptions{parskip=half}}
\makeatother
\usepackage{xcolor}
\usepackage{color}
\usepackage{fancyvrb}
\newcommand{\VerbBar}{|}
\newcommand{\VERB}{\Verb[commandchars=\\\{\}]}
\DefineVerbatimEnvironment{Highlighting}{Verbatim}{commandchars=\\\{\}}
% Add ',fontsize=\small' for more characters per line
\newenvironment{Shaded}{}{}
\newcommand{\AlertTok}[1]{\textcolor[rgb]{1.00,0.00,0.00}{\textbf{#1}}}
\newcommand{\AnnotationTok}[1]{\textcolor[rgb]{0.38,0.63,0.69}{\textbf{\textit{#1}}}}
\newcommand{\AttributeTok}[1]{\textcolor[rgb]{0.49,0.56,0.16}{#1}}
\newcommand{\BaseNTok}[1]{\textcolor[rgb]{0.25,0.63,0.44}{#1}}
\newcommand{\BuiltInTok}[1]{\textcolor[rgb]{0.00,0.50,0.00}{#1}}
\newcommand{\CharTok}[1]{\textcolor[rgb]{0.25,0.44,0.63}{#1}}
\newcommand{\CommentTok}[1]{\textcolor[rgb]{0.38,0.63,0.69}{\textit{#1}}}
\newcommand{\CommentVarTok}[1]{\textcolor[rgb]{0.38,0.63,0.69}{\textbf{\textit{#1}}}}
\newcommand{\ConstantTok}[1]{\textcolor[rgb]{0.53,0.00,0.00}{#1}}
\newcommand{\ControlFlowTok}[1]{\textcolor[rgb]{0.00,0.44,0.13}{\textbf{#1}}}
\newcommand{\DataTypeTok}[1]{\textcolor[rgb]{0.56,0.13,0.00}{#1}}
\newcommand{\DecValTok}[1]{\textcolor[rgb]{0.25,0.63,0.44}{#1}}
\newcommand{\DocumentationTok}[1]{\textcolor[rgb]{0.73,0.13,0.13}{\textit{#1}}}
\newcommand{\ErrorTok}[1]{\textcolor[rgb]{1.00,0.00,0.00}{\textbf{#1}}}
\newcommand{\ExtensionTok}[1]{#1}
\newcommand{\FloatTok}[1]{\textcolor[rgb]{0.25,0.63,0.44}{#1}}
\newcommand{\FunctionTok}[1]{\textcolor[rgb]{0.02,0.16,0.49}{#1}}
\newcommand{\ImportTok}[1]{\textcolor[rgb]{0.00,0.50,0.00}{\textbf{#1}}}
\newcommand{\InformationTok}[1]{\textcolor[rgb]{0.38,0.63,0.69}{\textbf{\textit{#1}}}}
\newcommand{\KeywordTok}[1]{\textcolor[rgb]{0.00,0.44,0.13}{\textbf{#1}}}
\newcommand{\NormalTok}[1]{#1}
\newcommand{\OperatorTok}[1]{\textcolor[rgb]{0.40,0.40,0.40}{#1}}
\newcommand{\OtherTok}[1]{\textcolor[rgb]{0.00,0.44,0.13}{#1}}
\newcommand{\PreprocessorTok}[1]{\textcolor[rgb]{0.74,0.48,0.00}{#1}}
\newcommand{\RegionMarkerTok}[1]{#1}
\newcommand{\SpecialCharTok}[1]{\textcolor[rgb]{0.25,0.44,0.63}{#1}}
\newcommand{\SpecialStringTok}[1]{\textcolor[rgb]{0.73,0.40,0.53}{#1}}
\newcommand{\StringTok}[1]{\textcolor[rgb]{0.25,0.44,0.63}{#1}}
\newcommand{\VariableTok}[1]{\textcolor[rgb]{0.10,0.09,0.49}{#1}}
\newcommand{\VerbatimStringTok}[1]{\textcolor[rgb]{0.25,0.44,0.63}{#1}}
\newcommand{\WarningTok}[1]{\textcolor[rgb]{0.38,0.63,0.69}{\textbf{\textit{#1}}}}
\setlength{\emergencystretch}{3em} % prevent overfull lines
\providecommand{\tightlist}{%
  \setlength{\itemsep}{0pt}\setlength{\parskip}{0pt}}
\setcounter{secnumdepth}{-\maxdimen} % remove section numbering
\ifLuaTeX
  \usepackage{selnolig}  % disable illegal ligatures
\fi
\IfFileExists{bookmark.sty}{\usepackage{bookmark}}{\usepackage{hyperref}}
\IfFileExists{xurl.sty}{\usepackage{xurl}}{} % add URL line breaks if available
\urlstyle{same}
\hypersetup{
  hidelinks,
  pdfcreator={LaTeX via pandoc}}

\author{}
\date{}

\begin{document}

\hypertarget{satspath-docker-cheat-sheet}{%
\section{🐳 SatsPath Docker Cheat
Sheet}\label{satspath-docker-cheat-sheet}}

Esta guía contiene todos los comandos necesarios para construir,
ejecutar y probar la funcionalidad completa de \textbf{SatsPath} usando
los contenedores de Docker de manera segura.

\begin{quote}
\textbf{NOTE:} Dado que el CLI se ejecuta de forma efímera, el formato
base para correr cualquier comando es:
\texttt{docker\ compose\ run\ -\/-rm\ satspath-cli\ \textless{}comando\textgreater{}}
\end{quote}

\begin{center}\rule{0.5\linewidth}{0.5pt}\end{center}

\hypertarget{construcciuxf3n-y-entorno-base}{%
\subsection{1. Construcción y Entorno
Base}\label{construcciuxf3n-y-entorno-base}}

Antes de ejecutar comandos, necesitas construir las imágenes e iniciar
el daemon en segundo plano.

\begin{Shaded}
\begin{Highlighting}[]
\CommentTok{\# Construir todas las imágenes (CLI, Daemon y Ark Bridge)}
\ExtensionTok{docker}\NormalTok{ compose build}

\CommentTok{\# Levantar el Daemon local (satspathd) en segundo plano}
\ExtensionTok{docker}\NormalTok{ compose up }\AttributeTok{{-}d}\NormalTok{ satspathd}

\CommentTok{\# Verificar que el Daemon esté corriendo y saludable (Healthcheck)}
\ExtensionTok{docker}\NormalTok{ compose ps}
\end{Highlighting}
\end{Shaded}

\begin{center}\rule{0.5\linewidth}{0.5pt}\end{center}

\hypertarget{inicializaciuxf3n-de-billetera-y-perfil}{%
\subsection{2. Inicialización de Billetera y
Perfil}\label{inicializaciuxf3n-de-billetera-y-perfil}}

SatsPath es una billetera ``receiver-profile''. Gestiona tu identidad y
tus métodos de pago públicos, sin custodiar fondos.

\begin{Shaded}
\begin{Highlighting}[]
\CommentTok{\# 1. Inicializar las llaves de identidad criptográfica}
\ExtensionTok{docker}\NormalTok{ compose run }\AttributeTok{{-}{-}rm}\NormalTok{ satspath{-}cli wallet init}

\CommentTok{\# 2. Registrar tu alias y configurar tus métodos de pago públicos}
\ExtensionTok{docker}\NormalTok{ compose run }\AttributeTok{{-}{-}rm}\NormalTok{ satspath{-}cli wallet add{-}methods bob@satspath.local }\DataTypeTok{\textbackslash{}}
    \AttributeTok{{-}{-}lightning{-}address}\NormalTok{ bob@getalby.com }\DataTypeTok{\textbackslash{}}
    \AttributeTok{{-}{-}onchain{-}address}\NormalTok{ tb1qw508d6qejxtdg4y5r3zarvary0c5xw7kxpjzsx }\DataTypeTok{\textbackslash{}}
    \AttributeTok{{-}{-}ark{-}server}\NormalTok{ https://ark{-}testnet.example.com }\DataTypeTok{\textbackslash{}}
    \AttributeTok{{-}{-}ark{-}pubkey}\NormalTok{ 025b90f4a...}

\CommentTok{\# 3. Ver tu perfil actual y el estado de la firma criptográfica}
\ExtensionTok{docker}\NormalTok{ compose run }\AttributeTok{{-}{-}rm}\NormalTok{ satspath{-}cli wallet show}
\end{Highlighting}
\end{Shaded}

\begin{center}\rule{0.5\linewidth}{0.5pt}\end{center}

\hypertarget{resolviendo-pagos-simulaciones-y-quotes}{%
\subsection{3. Resolviendo Pagos (Simulaciones y
Quotes)}\label{resolviendo-pagos-simulaciones-y-quotes}}

Antes de enviar un pago, el sistema de SatsPath debe evaluar las tarifas
de red, urgencia y métodos de la contraparte para decidir la ruta
óptima.

\begin{Shaded}
\begin{Highlighting}[]
\CommentTok{\# Consultar el perfil público de otro usuario}
\ExtensionTok{docker}\NormalTok{ compose run }\AttributeTok{{-}{-}rm}\NormalTok{ satspath{-}cli resolve bob@satspath.local}

\CommentTok{\# Generar un "Quote" de pago por 50,000 sats (SatsPath decide la mejor ruta)}
\ExtensionTok{docker}\NormalTok{ compose run }\AttributeTok{{-}{-}rm}\NormalTok{ satspath{-}cli quote bob@satspath.local 50000}

\CommentTok{\# Forzar una urgencia de pago rápida (altera la selección de red Lightning vs On{-}chain)}
\ExtensionTok{docker}\NormalTok{ compose run }\AttributeTok{{-}{-}rm}\NormalTok{ satspath{-}cli quote bob@satspath.local 250000 }\AttributeTok{{-}{-}urgency}\NormalTok{ fast}
\end{Highlighting}
\end{Shaded}

\begin{center}\rule{0.5\linewidth}{0.5pt}\end{center}

\hypertarget{ejecuciuxf3n-de-pagos-y-swaps}{%
\subsection{4. Ejecución de Pagos y
Swaps}\label{ejecuciuxf3n-de-pagos-y-swaps}}

El comando \texttt{pay} muestra la instrucción de pago. Si le agregas
los flags experimentales, puedes simular la ejecución a través del
\texttt{LightningExecutor} u orquestar Swaps Submarinos.

\begin{Shaded}
\begin{Highlighting}[]
\CommentTok{\# Pagar a un destinatario y ver las instrucciones del QR (No mueve fondos)}
\ExtensionTok{docker}\NormalTok{ compose run }\AttributeTok{{-}{-}rm}\NormalTok{ satspath{-}cli pay bob@satspath.local 1000}

\CommentTok{\# Simular la ejecución experimental del pago/swap en testnet}
\ExtensionTok{docker}\NormalTok{ compose run }\AttributeTok{{-}{-}rm}\NormalTok{ satspath{-}cli pay bob@satspath.local 100000 }\AttributeTok{{-}{-}testnet} \AttributeTok{{-}{-}experimental{-}swaps}
\end{Highlighting}
\end{Shaded}

\begin{center}\rule{0.5\linewidth}{0.5pt}\end{center}

\hypertarget{el-flujo-de-invitaciones-invites-claims}{%
\subsection{5. El Flujo de Invitaciones (Invites \&
Claims)}\label{el-flujo-de-invitaciones-invites-claims}}

Si le pagas a un usuario que \textbf{no} tiene un perfil registrado,
SatsPath genera automáticamente un enlace de invitación.

\begin{Shaded}
\begin{Highlighting}[]
\CommentTok{\# 1. Intentar pagarle a un usuario nuevo (Genera la invitación)}
\ExtensionTok{docker}\NormalTok{ compose run }\AttributeTok{{-}{-}rm}\NormalTok{ satspath{-}cli pay nuevo\_usuario@satspath.local 50000}

\CommentTok{\# El comando anterior te devolverá un enlace (ej: https://satspath.local/claim?alias\_hash=...)}
\CommentTok{\# 2. El nuevo usuario reclama su perfil usando el enlace}
\ExtensionTok{docker}\NormalTok{ compose run }\AttributeTok{{-}{-}rm}\NormalTok{ satspath{-}cli claim }\StringTok{"https://satspath.local/claim?alias\_hash=abc123def456\&amount=50000"}
\end{Highlighting}
\end{Shaded}

\begin{center}\rule{0.5\linewidth}{0.5pt}\end{center}

\hypertarget{mantenimiento-de-seguridad-rotaciuxf3n-de-llaves}{%
\subsection{6. Mantenimiento de Seguridad (Rotación de
Llaves)}\label{mantenimiento-de-seguridad-rotaciuxf3n-de-llaves}}

Si crees que tu llave de identidad fue comprometida, debes rotarla.
SatsPath usa rotación criptográfica con pruebas \texttt{KeyRotation}.

\begin{Shaded}
\begin{Highlighting}[]
\CommentTok{\# 1. Rotar las llaves de identidad (crea una nueva llave y adjunta la prueba)}
\ExtensionTok{docker}\NormalTok{ compose run }\AttributeTok{{-}{-}rm}\NormalTok{ satspath{-}cli wallet rotate}

\CommentTok{\# 2. Validar el estado del perfil}
\CommentTok{\# Deberías ver un aviso indicando que el perfil ha sido rotado recientemente}
\ExtensionTok{docker}\NormalTok{ compose run }\AttributeTok{{-}{-}rm}\NormalTok{ satspath{-}cli wallet show}
\end{Highlighting}
\end{Shaded}

\begin{center}\rule{0.5\linewidth}{0.5pt}\end{center}

\hypertarget{ark-bridge-opcional}{%
\subsection{7. ARK Bridge (Opcional)}\label{ark-bridge-opcional}}

Si necesitas utilizar el puente TypeScript de Ark para validación
avanzada del cliente, debes levantar el perfil \texttt{bridge}.

\begin{Shaded}
\begin{Highlighting}[]
\CommentTok{\# Iniciar el contenedor de ark{-}bridge}
\ExtensionTok{docker}\NormalTok{ compose }\AttributeTok{{-}{-}profile}\NormalTok{ bridge up }\AttributeTok{{-}d}\NormalTok{ ark{-}bridge}

\CommentTok{\# Detener todos los servicios al terminar}
\ExtensionTok{docker}\NormalTok{ compose }\AttributeTok{{-}{-}profile}\NormalTok{ bridge down}
\end{Highlighting}
\end{Shaded}


\end{document}
